\chapter{Bridge Moment Hierarchy and the Second-Order Tail Theorem}
\label{chap:g024-bridge-moment-tail}

\status{Exact SOH-G024 bridge-moment hierarchy and analytic proof of $H_y''(q)>0$ for $q\ge1/4$ and $0<|y|<1/2$. Historically the compact core $0\le q<1/4$ remained open here; Chapters~53--54 subsequently close it. The current unresolved derivative frontier is $m\ge3$.}

\section{Bridge score hierarchy}
With
\[
d\mu_u(r)=\frac{r^2K(u+r)K(u-r)}{C(u)}\,dr,
\qquad L=-\log K,
\]
define
\[
D_u(r)=L'(u+r)-L'(u-r).
\]
The logarithmic density derivative is
\[
\partial_r\log\mu_u(r)=\frac2r-D_u(r).
\]
Integration by parts against $r^{2n+1}$ gives
\[
\boxed{\mathbb E[r^{2n+1}D_u(r)]=(2n+3)\mathbb E[r^{2n}]},
\]
and in particular
\[
\boxed{\mathbb E[rD_u(r)]=3.}
\]

\section{Conservative moment concentration}
At this stage the G004 curvature margin gives $L''>10$, hence $D_u(r)>20r$ for $r>0$. Therefore
\[
\boxed{\mathbb E[r^{2n+2}]<\frac{2n+3}{20}\mathbb E[r^{2n}]},
\]
so
\[
\boxed{\mathbb E[r^{2n}]<\frac{(2n+1)!!}{20^n}},
\qquad
\boxed{\mathbb E[r^2]<\frac3{20}},
\]
and, for $0\le\lambda<10$,
\[
\boxed{\mathbb E[e^{\lambda r^2}]\le(1-\lambda/10)^{-3/2}}.
\]
Chapter~53 sharpens the denominator pair $(20,10)$ to $(34,17)$ using $L''>17$.

\begin{figure}[htbp]
\centering
\includegraphics[width=0.76\textwidth]{bridge_moment_envelope.png}
\caption{The sharpened bridge-moment envelope proved in Chapter~53, $(2n+1)!!/34^n$. It is displayed here to connect the conservative hierarchy of the present chapter with the later strengthened result.}
\label{fig:bridge-envelope-ch52}
\end{figure}

\section{Companion upper curvature envelope}
For the theta channels used in G004, the channel curvature and the same ratio bounds give
\[
\boxed{10<L''(s)<21e^{2|s|}}.
\]
Hence
\[
B_u(r)=L''(u+r)+L''(u-r)<42e^{2u+2|r|}.
\]
Using
\[
2|r|\le\frac52r^2+\frac25
\]
and the conservative MGF yields
\[
\boxed{\mathbb E[B_u]<36e\,e^{2u}}.
\]

\section{First-order surplus in the tail}
The lower channel-curvature envelope gives
\[
R(u)>4\pi(e^{2u}-1)-4u.
\]
With $T_y(u)=2|y|\tanh(2|y|u)<u$,
\[
N_y(u)>12(e^{2u}-1)-5u.
\]
For $u\ge1/2$ this implies
\[
\boxed{N_y(u)>6e^{2u}}.
\]
Therefore
\[
N_y(u)^2>36e^{4u}>\mathbb E[B_u]
\]
in that region.

\section{SOH-G024 second-order tail theorem}
Insert the preceding estimates into the exact bridge normal form of Chapter~51. Since the variance and tilt-curvature terms are non-negative,
\[
\boxed{H_y''(q)>0\qquad\left(q\ge\frac14,\ 0<|y|<\frac12\right).}
\]

\begin{theorem}[SOH-G024 second-order tail theorem]
For the Riemann-kernel external Dimitrov--Xu family, the second complete-monotonicity inequality holds strictly throughout the tail region $q\ge1/4$.
\end{theorem}

\section{Historical compact core and later closure}
At this point the second-order obligation had been reduced to
\[
\boxed{0\le q<\frac14}.
\]
Chapter~53 shrinks this historical core to $0\le q<1/9$, and Chapter~54 closes it globally. The conservative results of this chapter remain useful as independent subtheorems and as checks on the later sharpened estimates.

\gap{The former second-order core is no longer a current open problem; Chapter~54 proves global second-order positivity. The present unresolved G024 frontier is $m\ge3$, together with strict external Fourier positivity, SOH-G003, SOH-C005, PF$_3$, PF$_\infty$, and RH.}
